% part: intuitionistic-logic
% chapter: semantics
% section: topological-semantics

\documentclass[../../../include/open-logic-section]{subfiles}

\begin{document}

\olfileid[pt]{int}{sem}{top}

\olsection{Semântica topológica}

Outra maneira de fornecer uma semântica para a lógica intuicionista é
usar o conceito matemático de topologia.

\begin{defn}
  Seja $X$ um conjunto. Uma \emph{topologia sobre~$X$} é um conjunto
  $\Top{O} \subseteq \Pow{X}$ que satisfaz as propriedades abaixo. Os
  !!{element}s de~$\Top{O}$ são chamados de \emph{conjuntos abertos}
  da topologia. O conjunto $X$, juntamente com $\Top{O}$, é chamado de
  \emph{espaço topológico}.
  \begin{enumerate}
  \item O conjunto vazio e o espaço inteiro são abertos: $\emptyset$, $X \in
    \Top{O}$.
  \item Os conjuntos abertos são fechados sob interseções finitas: se
    $U$, $V \in \Top{O}$, então $U \cap V \in \Top{O}$.
  \item Os conjuntos abertos são fechados sob uniões arbitrárias: se
    $U_i \in \Top{O}$ para todo $i \in I$, então $\bigcup \Setabs{U_i}{i \in
      I} \in \Top{O}$.
  \end{enumerate}
\end{defn}

Podemos escrever $X$ para uma topologia quando a coleção de conjuntos
abertos puder ser inferida do contexto; observe, ainda assim, que $X$
só pode ser chamado de topologia depois de ser dotado de conjuntos
abertos.

\begin{defn}
  Um \emph{modelo topológico} da lógica proposicional intuicionista é
  uma tripla $\mModel{X} = \tuple{X, \Top{O}, V}$, em que $\Top{O}$ é
  uma topologia sobre~$X$ e $V$ é uma função que atribui a cada
  variável proposicional um conjunto aberto de $\Top{O}$.

  Dado um modelo topológico~$\mModel{X}$, podemos definir
  $\Prop{X}{!A}$ indutivamente da seguinte maneira:
  \begin{enumerate}
  \item $\Prop{X}{\lfalse} = \emptyset$
  \item $\Prop{X}{p} = V(p)$
  \item $\Prop{X}{!A \land !B} = \Prop{X}{!A} \cap \Prop{X}{!B}$
  \item $\Prop{X}{!A \lor !B} = \Prop{X}{!A} \cup \Prop{X}{!B}$
  \item $\Prop{X}{!A \lif !B} = \Interior{(X \setminus \Prop{X}{!A}) \cup
    \Prop{X}{!B}}$
  \end{enumerate}
  Aqui, $\Interior{V}$ é a função que leva um conjunto $V \subseteq X$
  a seu \emph{interior}, isto é, à união de todos os conjuntos abertos
  que ele contém. Em outras palavras,
  \[
  \Interior{V} = \bigcup \Setabs{U}{U \subseteq V \text{ e } U \in \Top{O}}.
  \]
\end{defn}

Observe que o interior de qualquer conjunto é sempre aberto, pois é
uma união de conjuntos abertos. Portanto, $\Prop{X}{!A}$ é sempre um
conjunto aberto.

Embora a semântica topológica seja muito abstrata, há maneiras de
entendê-la que podem motivá-la. Suponha que os !!{element}s, ou
``pontos'', de $X$ sejam pontos nos quais afirmações podem ser
avaliadas. O conjunto de todos os pontos nos quais $!A$ é verdadeira é
a proposição expressa por~$!A$. Nem todo conjunto de pontos é uma
proposição possível; somente os !!{element}s de $\Top{O}$ o são.
$!A \Entails !B$ sse $!B$ é verdadeira em todo ponto no qual~$!A$ é
verdadeira, isto é, $\Prop{X}{!A} \subseteq \Prop{X}{!B}$, para
todo~$X$. A afirmação absurda~$\lfalse$ nunca é verdadeira; portanto,
$\Prop{X}{\lfalse} = \emptyset$.

Como devem as proposições expressas por $!B \land !C$, $!B \lor !C$
e $!B \lif !C$ relacionar-se com as expressas por $!B$ e~$!C$ para que
valham as leis intuicionisticamente válidas, isto é, para que
$!A \Proves !B$ sse $\Prop{X}{!A} \subset \Prop{X}{!B}$? Exigimos
$\lfalse \Proves !A$ para todo $!A$, o que é satisfeito porque
$\emptyset \subseteq U$ para todo $U$. Como $!B \land !C \Proves !B$,
exigimos $\Prop{X}{!B \land !C} \subseteq \Prop{X}{!B}$ e, de modo
semelhante, $\Prop{X}{!B \land !C} \subseteq \Prop{X}{!C}$. O maior
conjunto que satisfaz $W \subseteq U$ e $W \subseteq V$ é $U \cap V$.
Reciprocamente, $!B \Proves !B \lor !C$ e $!C \Proves !B \lor !C$;
portanto, exigimos $\Prop{X}{!B} \subseteq \Prop{X}{!B \lor !C}$ e
$\Prop{X}{!C} \subseteq \Prop{X}{!B \lor !C}$. O menor conjunto~$W$
tal que $U \subseteq W$ e $V \subseteq W$ é $U \cup V$.

A definição de $\lif$ é delicada: $!A \lif !B$ expressa a proposição
mais fraca que, combinada com $!A$, implica $!B$. Que $!A \lif !B$,
combinada com $!A$, implica~$!B$ fica claro por
$(!A \lif !B) \land !A \Proves !B $. Assim, $\Prop{X}{!A \lif !B}$
deve ser o maior conjunto aberto tal que
$\Prop{X}{!A \lif !B} \cap \Prop{X}{!A} \subset \Prop{X}{!B}$, o que
leva à nossa definição.

\end{document}
